\documentclass[UTF8,11pt,a4paper,fontset=fandol]{ctexart}

\usepackage[a4paper,margin=27mm]{geometry}
\usepackage{amsmath,amssymb,amsthm,mathtools}
\usepackage{booktabs}
\usepackage{array}
\usepackage{enumitem}
\usepackage{xcolor}
\usepackage{xurl}
\usepackage{hyperref}
\usepackage[nameinlink,noabbrev]{cleveref}

\hypersetup{
  unicode=true,
  colorlinks=true,
  linkcolor=blue!55!black,
  citecolor=blue!55!black,
  urlcolor=blue!55!black,
  pdftitle={饱和 6-Sperner 数与 7-Sperner 数的精确值},
  pdfauthor={Jiazhi Chen},
  pdfsubject={极值集合论与饱和 Sperner 系统},
  pdfkeywords={饱和 Sperner 族, 布尔格, blocker, 极值集合论, SAT 证书, Lean 4},
  pdflang={zh-CN},
  pdfdisplaydoctitle=true,
  bookmarksnumbered=true
}

\newtheorem{theorem}{定理}[section]
\newtheorem{proposition}[theorem]{命题}
\newtheorem{lemma}[theorem]{引理}
\newtheorem{corollary}[theorem]{推论}
\theoremstyle{definition}
\newtheorem{definition}[theorem]{定义}
\theoremstyle{remark}
\newtheorem{remark}[theorem]{注}
\renewcommand{\proofname}{证明}

\crefformat{section}{第#2#1#3节}
\crefformat{table}{表#2#1#3}
\crefformat{theorem}{定理#2#1#3}
\crefformat{proposition}{命题#2#1#3}
\crefformat{lemma}{引理#2#1#3}
\crefformat{corollary}{推论#2#1#3}

\newcommand{\powerset}{\mathcal{P}}
\newcommand{\family}{\mathcal{F}}
\newcommand{\smallcores}{\mathcal{S}}
\newcommand{\largecores}{\mathcal{L}}
\newcommand{\sat}{\operatorname{sat}}
\newcommand{\blocker}{\mathcal{B}}
\newcommand{\bits}[1]{\mathtt{#1}}

\title{饱和 6-Sperner 数与 7-Sperner 数的精确值}
\author{Jiazhi Chen\\
\small 浙江省杭州第二中学\\
\small \href{mailto:lailai0x394@gmail.com}{lailai0x394@gmail.com}\quad
\href{https://orcid.org/0009-0008-9790-9881}{ORCID: 0009-0008-9790-9881}}
\date{}

\begin{document}

\pagestyle{plain}
\maketitle

\begin{abstract}
最终稳定的饱和 6-Sperner 数与 7-Sperner 数分别为 30 和 55。下界从有限
典范归约出发。齐次原子出现后，相邻层会在任意有限迹底集上产生互 blocker
pair。精确值 $m(2,3)=9$、$m(2,4)=12$ 与 $m(3,3)=14$ 解决六层情形，
并给出七层情形的初始轮廓。$m(3,3)$ 的等号情形是 Fano 平面。对相邻层及其余
中间 blocker pair 的分类排除了稳定范围内所有不超过 54 元的族。匹配的 55 元构造使用
八元素核心和一个与之不交的齐次原子。它包含 28 个小成员和 27 个大成员，故
复合还给出
\[
  \sat(5j+2+s)\leq 2^s(28^j+27^j)
  \qquad(j\geq1,\ 0\leq s\leq4).
\]
任意底集上的论证及两个精确值均已在 Lean 4 中形式化。有限分类通过已证明的
归约连接到任意有限 blocker pair。
\end{abstract}

\noindent\textbf{关键词：}饱和 Sperner 族；布尔格；blocker；极值集合论；
SAT 证书；Lean 4

\section{引言}

设 $X$ 为有限集。若族 $\family\subseteq\powerset(X)$ 不含长度为 $k+1$ 的
严格链，则称它为 $k$-Sperner 族。若加入任何缺失的子集都会产生这样的链，
则称它是饱和的。当 $|X|=n$ 时，相应最小基数记为 $\sat(n,k)$。禁链条件是
局部的，饱和性却要求每个缺失集合在加入后都能嵌入一条禁链。

Gerbner 等人~\cite{gerbner2013} 证明，对固定的 $k$，当底集足够大时，该
最小值会稳定。其最终稳定值记为 $\sat(k)$。Morrison、Noel 和
Scott~\cite{mns2014} 从饱和反链建立分层构造，并给出一个 30 元的六层构造。
他们的复合运算还能把一个小构造转化为无穷多个参数的上界。

Martin 和 Veldt~\cite{martinveldt2025} 沿用这一框架，构造了一个 56 元饱和
7-Sperner 族，并询问它是否最优。该构造在参数 $5j+2+s$ 处给出上界
$2^{s+1}28^j$。后来的机器学习搜索找到一个 108 元饱和 8-Sperner
族~\cite{patternboost}，但没有解决七层问题。

六层与七层的精确值分别为 30 和 55。共同起点是有限典范归约。相邻典范层产生
互 blocker pair，因而稳定下界转化为 clutter 的局部极值问题。
$m(2,3)=9$、$m(2,4)=12$ 与 $m(3,3)=14$ 解决六层情形，并把七层轮廓压缩到
少数可能。$m(3,3)$ 的等号情形由 Fano 平面控制。相邻层分析排除最小轮廓，
剩余的 54 元情形归结为四种中间拆分。

匹配构造由八元素核心和公共齐次原子上的七个饱和反链组成。每一层由极大避集
规则验证，前驱表记录分层关系。构造包含 28 个小成员和 27 个大成员，因此每次
有限复合都严格改进原有的 $28/28$ 因子。

部分中间 blocker 分类虽是有限的，却不适合在正文逐项展开。Supplement 证明
从任意有限 clutter 到这些有限域的归约，并给出一条完整示例。精确值论证也已
在 Lean 4 中形式化。直接验证器和 DRAT 证明分别复核构造及七点核心模板类。

\section{预备知识}
\label{sec:preliminaries}

我们用 $A\subsetneq B$ 表示严格包含。长度为 $r$ 的\emph{严格链}是一个序列
$A_0\subsetneq\cdots\subsetneq A_{r-1}$。

\begin{definition}
若族 $\family\subseteq\powerset(X)$ 不含长度为 $k+1$ 的严格链，则称其为
\emph{$k$-Sperner 族}。如果它是 $k$-Sperner 族，并且对每个
$S\in\powerset(X)\setminus\family$，族 $\family\cup\{S\}$ 都含有长度为
$k+1$ 的严格链，则称其为\emph{饱和 $k$-Sperner 族}。
\end{definition}

饱和 1-Sperner 族恰好就是极大反链。下面的概念用于组织若干这样的反链。

\begin{definition}
称两两不交的族 $\mathcal A_0,\ldots,\mathcal A_{k-1}$ 是\emph{分层的}，
如果对每个 $i>0$ 以及每个 $B\in\mathcal A_i$，都存在
$A\in\mathcal A_{i-1}$ 使得 $A\subsetneq B$。
\end{definition}

我们使用 Morrison--Noel--Scott 的下述定理，并采用 Martin--Veldt 使用的表述。

\begin{theorem}[分层饱和反链~\cite{mns2014,martinveldt2025}]
\label{thm:layered}
如果 $\mathcal A_0,\ldots,\mathcal A_{k-1}$ 两两不交，并且构成分层的饱和
反链，那么它们的并是一个饱和 $k$-Sperner 族。
\end{theorem}

我们的各层使用一个公共块。固定有限核心 $C$ 和一个与 $C$ 不交的非空集合
$H$。\emph{小模板}是子集 $S\subseteq C$；\emph{大模板}形如 $L\cup H$，
其中 $L\subseteq C$。若核心集合 $L$ 不包含
$\smallcores\subseteq\powerset(C)$ 的任何成员，则称 $L$ 是
\emph{$\smallcores$-自由的}。

按照 Morrison--Noel--Scott 的定义，族的一个\emph{原子}是满足下述性质的
包含极大集合 $H$：族中的每个成员要么包含 $H$，要么与 $H$ 不交。当
$|H|>2$ 时，称该原子是\emph{齐次的}。

有限族 $\family$ 的\emph{典范层}由逐次删除包含极小成员得到。具体地，令
$\mathcal R_0=\family$，令 $\mathcal A_i$ 为 $\mathcal R_i$ 的全部包含极小
成员，并令 $\mathcal R_{i+1}=\mathcal R_i\setminus\mathcal A_i$。

\begin{lemma}[极大无小核证书]
\label{lem:maximal-free}
设 $\smallcores\subseteq\powerset(C)$ 是一个反链，$\largecores$ 是 $C$ 中
所有包含极大的 $\smallcores$-自由子集所成的集合。那么
\[
  \mathcal A(\smallcores,\largecores)
  =\smallcores\ \cup\ \{L\cup H:L\in\largecores\}
\]
是 $\powerset(C\cup H)$ 中的一个饱和反链。
\end{lemma}

\begin{proof}
各个小模板彼此不可比。不同的极大 $\smallcores$-自由集合彼此不可比，所以各个
大模板也彼此不可比。若 $L\in\largecores$，小模板 $S$ 不可能包含于
$L\cup H$，否则 $S\subseteq L$，这与 $L$ 是 $\smallcores$-自由的相矛盾。
大模板也不可能包含于小模板，因为大模板含有 $H$。因此 $\mathcal A$ 是反链。

现取 $\mathcal A$ 外部的 $T\subseteq C\cup H$，并令 $R=T\cap C$。如果
$R$ 包含某个 $S\in\smallcores$，那么 $S\subsetneq T$。否则，$R$ 是
$\smallcores$-自由的。由于 $C$ 有限，$R$ 包含于某个包含极大的
$\smallcores$-自由集合 $L\in\largecores$，从而 $T\subseteq L\cup H$。
这个包含是严格的，因为 $T$ 不属于 $\mathcal A$。所以每个外部集合都与
$\mathcal A$ 的某个成员严格可比，饱和性得证。
\end{proof}

\section{Blocker 与稳定下界}
\label{sec:blockers}

设 $U$ 为有限集。若集合与族 $\mathcal H\subseteq\powerset(U)$ 的每个成员
相交，则称其为 $\mathcal H$ 的\emph{横截集}。所有包含极小横截集组成的族
称为 $\mathcal H$ 的 \emph{blocker}，记为 $\blocker(\mathcal H)$。若族中
任意两个不同成员互不包含，则称该族为 \emph{clutter}。对正整数 $p,q$，定义
\[
  m(p,q)=\min\bigl\{|\mathcal S|+|\mathcal C|:
  \mathcal S=\blocker(\mathcal C),\
  \mathcal C=\blocker(\mathcal S),\
  \min_{S\in\mathcal S}|S|\geq p,\
  \min_{C\in\mathcal C}|C|\geq q\bigr\},
\]
其中最小值遍历所有有限底集及满足条件的 clutter。

\begin{lemma}[Blocker 对合与私有见证]
\label{lem:blocker-basic}
对每个有限 clutter $\mathcal H$，有
$\blocker(\blocker(\mathcal H))=\mathcal H$。此外，若
$T\in\blocker(\mathcal H)$ 且 $x\in T$，则存在
$E_x\in\mathcal H$ 使得 $T\cap E_x=\{x\}$。因此
$|T|\leq|\mathcal H|$。
\end{lemma}

\begin{proof}
$T$ 的极小性说明 $T\setminus\{x\}$ 会漏掉某个 $E_x$，但 $T$ 与
$E_x$ 相交。因此 $T\cap E_x=\{x\}$。不同点对应不同私有见证，故得到基数
不等式。

每个 $E\in\mathcal H$ 都是 $\blocker(\mathcal H)$ 的横截集。若其真子集
$D$ 也是横截集，则 clutter 性说明 $D$ 不包含 $\mathcal H$ 的成员。于是
$U\setminus D$ 与 $\mathcal H$ 的每个成员相交，并包含一个与 $D$ 不交的极小
横截集，矛盾。因此
$\mathcal H\subseteq\blocker(\blocker(\mathcal H))$。反向包含中，后一
blocker 的每个成员都包含 $\mathcal H$ 的某个成员，极小性迫使二者相等。
\end{proof}

\begin{lemma}[分层反链的向上步进]
\label{lem:upward-step}
设 $\mathcal A_0,\ldots,\mathcal A_{k-1}$ 是两两不交、满足分层条件的
饱和反链。若 $D\in\mathcal A_i$ 且 $i<k-1$，则存在
$E\in\mathcal A_{i+1}$ 使 $D\subsetneq E$。因此，从 $D$ 出发可得到一条
严格链，它在第 $i$ 层至第 $k-1$ 层各取一个成员。
\end{lemma}

\begin{proof}
两层不交给出 $D\notin\mathcal A_{i+1}$。由 $\mathcal A_{i+1}$ 的饱和性，
存在 $E\in\mathcal A_{i+1}$ 与 $D$ 严格可比。若 $E\subsetneq D$，分层条件
给出 $C\in\mathcal A_i$ 使 $C\subsetneq E$。于是
$C\subsetneq D$，与 $\mathcal A_i$ 为反链矛盾。因此
$D\subsetneq E$。逐层重复该论证即可得到所需严格链。
\end{proof}

下面的定理给出两个全局下界实际使用的精确接口，并明确保证某个关联类至少含
三个点所需的数值阈值。

\begin{theorem}[有限典范归约]
\label{thm:canonical-reduction}
设 $k\geq2$，且 $\family\subseteq\powerset(X)$ 是有限底集 $X$ 上的饱和
$k$-Sperner 族。若
\[
  |\family|<2^{|X|},\qquad 3\cdot2^{|\family|}\leq|X|,       \tag{CR1}
\]
则 $\family$ 有一个齐次原子 $H$，并恰有 $k$ 个非空典范层
$\mathcal A_0,\ldots,\mathcal A_{k-1}$。这些层两两不交，是分层的饱和反链，
且其并为 $\family$。

令 $U=X\setminus H$。对 $0\leq i<k$，定义
\[
 \mathcal S_i=\{A\in\mathcal A_i:A\cap H=\varnothing\},
\]
\[
 \mathcal C_i=
 \{U\setminus(A\setminus H):A\in\mathcal A_i,\ H\subseteq A\}.
\]
对每个内部指标 $1\leq i\leq k-2$，有
\[
 \mathcal S_i=\blocker(\mathcal C_i),\qquad
 \mathcal C_i=\blocker(\mathcal S_i),                       \tag{CR2}
\]
$\mathcal S_i$ 的每个成员大小至少为 $i$，$\mathcal C_i$ 的每个成员大小
至少为 $k-1-i$，并且
\[
 |\mathcal A_i|=|\mathcal S_i|+|\mathcal C_i|.              \tag{CR3}
\]
\end{theorem}

\begin{proof}
给每个点 $x\in X$ 赋关联签名
$\sigma(x)=\{F\in\family:x\in F\}$。签名至多有 $2^{|\family|}$ 种。
由 (CR1) 和抽屉原理，某个签名类 $H$ 至少有三个点。族中每个成员要么包含整个
$H$，要么与 $H$ 不交；任何类外点都有不同签名，故都不能加入 $H$ 而保持这一
性质。因此 $H$ 是齐次原子。

逐次取极小成员使典范层两两不交、各自成反链并满足分层条件。$k$-Sperner 性
说明第 $k$ 层以上没有成员，所以前 $k$ 层的并为 $\family$。由 (CR1) 存在
外部集合 $T$。饱和性给出一条经过 $T$ 的 $(k+1)$ 项严格链；删去 $T$ 后得到
$\family$ 中的一条 $k$ 项严格链，它在每个典范层恰有一个成员。因此这 $k$ 层
均非空。

还需证明每层饱和。固定 $i$ 与 $D\notin\mathcal A_i$。取不同的
$x,y\in H$，令 $D'=(D\setminus H)\cup\{x\}$。集合 $D'$ 与 $H$ 相交但不
包含 $H$，故不属于 $\family$。取经过 $D'$ 的严格链并删去 $D'$；其中第 $i$
层成员 $A$ 与 $D'$ 可比。若 $A\subsetneq D'$，原子性质迫使
$A\cap H=\varnothing$，从而 $A\subseteq D$；等号会给出
$D\in\mathcal A_i$，故包含严格。若 $D'\subsetneq A$，则 $H\subseteq A$，
从而 $D\subseteq A$，且等号同样不可能。因此每个层外集合都与该层某成员严格
可比，$\mathcal A_i$ 是饱和反链。

固定内部指标 $i$，并取非空真子集 $Q\subset H$。对每个 $R\subseteq U$，
混合集 $R\cup Q$ 不属于 $\mathcal A_i$。层的饱和性与原子性质给出二择一：
存在 $S\in\mathcal S_i$ 使 $S\subseteq R$，或者存在大成员
$A\in\mathcal A_i$ 使 $R\subseteq A\setminus H$。反链性质说明每个
$S\in\mathcal S_i$ 与每个 $C\in\mathcal C_i$ 相交。若 $T$ 是
$\mathcal C_i$ 的横截集，它不包含于任何大迹，故第一种情形给出某个
$S\subseteq T$。于是 $\mathcal C_i$ 的极小横截集恰为 $\mathcal S_i$；
blocker 对合给出 (CR2) 的反向等式。

最后，$\mathcal A_i$ 的成员有一条含 $i$ 个严格前驱的链，所以小迹大小至少
为 $i$。由\cref{lem:upward-step}，每个大成员可向上延伸到最后一层。对
$k-1-i$ 次严格扩张各取一个新增点。这些点两两不同；初始大成员已经包含整个
$H$，所以新增点均在 $H$ 外，也均在初始核心迹之外。故相应补集大小至少为
$k-1-i$。原子把该层划分为小、大两部分，取补是单射，故 (CR3) 成立。
\end{proof}

\begin{theorem}
\label{thm:m23}
精确局部 blocker 参数为 $m(2,3)=9$。
\end{theorem}

\begin{proof}
取一个满足条件的 pair，并记
$a=|\mathcal S|$、$b=|\mathcal C|$。在两个方向使用私有见证不等式，得到
$a\geq3$ 与 $b\geq2$。

若 $b=2$，则 $\mathcal C$ 的两个成员不交，否则交中的点会给出单点 blocker。
从二者各取一点，至少得到 $3^2=9$ 个 $\mathcal S$ 成员。

设 $b=3$，三行为 $C_1,C_2,C_3$。三重交为空。对
$\{i,j,k\}=\{1,2,3\}$，令
\[
  A_i=C_i\setminus(C_j\cup C_k),\qquad
  P_i=(C_j\cap C_k)\setminus C_i,
\]
并记 $\alpha_i=|A_i|$、$\pi_i=|P_i|$。二点极小横截集要么从
$A_i,P_i$ 各取一点，要么从两个不同的 $P$ 类各取一点。三点极小横截集恰好
从三个 $A_i$ 各取一点。\cref{lem:blocker-basic}中的私有见证排除了其他
类型。因此
\[
  a=\alpha_1\alpha_2\alpha_3+
  \sum_{i=1}^3\alpha_i\pi_i+
  \pi_1\pi_2+\pi_1\pi_3+\pi_2\pi_3.
  \tag{1}
\]
行大小给出 $\alpha_i+\pi_j+\pi_k\geq3$。若某个 $\pi_i\geq2$，式 (1)
中包含 $\pi_i$ 的项至少贡献六。等号会迫使 $\pi_i=2$、对应行大小为三，且
剩余项全部为零；另外两条行不等式随后迫使三个 $\alpha$ 类都非空，矛盾。
故 $a\geq7$。若所有 $\pi_i\leq1$，利用式 (1) 对 $\alpha_i$ 的单调性，
在对称意义下仅需检查
\[
\begin{array}{ccc}
\toprule
(\pi_1,\pi_2,\pi_3)&
\min(\alpha_1,\alpha_2,\alpha_3)&\min a\\
\midrule
(0,0,0)&(3,3,3)&27\\
(0,0,1)&(2,2,3)&15\\
(0,1,1)&(1,2,2)&9\\
(1,1,1)&(1,1,1)&7\\
\bottomrule
\end{array}
\]
所以仍有 $a\geq7$。

以下设 $b\geq4$。若 $a=3$，则 $\mathcal S$ 的三行两两不交；否则一个
交点加上第三行中的一点会构成至多二点的横截集。逐行取一点给出
$b\geq2^3=8$。若 $a=4$，在 $\mathcal S$ 的四行上构造交图。两条顶点不交
的边会给出至多二点的横截集，所以交图匹配数至多为一。若有两条孤立行，
blocker 按不交分量分解，至少给出 $2\cdot2\cdot2=8$ 个横截集。若恰有一条
孤立行，连通三行分量至少有三个 blocker，故总数至少为 $2\cdot3=6$。若没有
孤立行，交图只能是三叶星。三条叶行两两不交；从各叶行取一点，至少得到
\[
  n_1n_2n_3-(n_1-1)(n_2-1)(n_3-1)\geq7
\]
个极小横截集，其中 $n_i\geq2$ 是叶行大小。因此 $b\geq6$。若 $a\geq5$，
总数至少为九。所有情形都给出 $a+b\geq9$。

为证明取等，取不同点 $x_1,x_2,y_1,y_2,w,z$，并令
\[
  \mathcal S=\{\{x_i,y_j\}:i,j\in\{1,2\}\}\cup\{\{w,z\}\},
\]
\[
  \mathcal C=\{\{x_1,x_2,w\},\{x_1,x_2,z\},
  \{y_1,y_2,w\},\{y_1,y_2,z\}\}.
\]
$\mathcal C$ 的极小横截集要么各取一个 $x_i,y_j$，要么等于 $\{w,z\}$。
所以 $\blocker(\mathcal C)=\mathcal S$。Blocker 对合给出
$\blocker(\mathcal S)=\mathcal C$，两侧总大小为九。
\end{proof}

\begin{theorem}
\label{thm:sat6}
最终稳定的饱和 6-Sperner 数为
\[
  \sat(6)=30.
\]
\end{theorem}

\begin{proof}
Morrison--Noel--Scott 对每个至少含八个点的底集构造了一个 30 元饱和
6-Sperner 族~\cite[Proposition~20]{mns2014}，故 $\sat(6)\leq30$。

在稳定区间内选取满足 $n\geq3\cdot2^{30}$ 的底集，并在其上选取极小族
$\family$。由上界构造，$|\family|\leq30$，所以
$3\cdot2^{|\family|}\leq3\cdot2^{30}\leq n$，且
$|\family|<2^n$。因此\cref{thm:canonical-reduction}给出六个非空典范层
$\mathcal A_0,\ldots,\mathcal A_5$。端点层非空；对指标一和四的迹 blocker
pair 使用私有见证，得到
\[
  |\mathcal A_0|,|\mathcal A_5|\geq1,\qquad
  |\mathcal A_1|,|\mathcal A_4|\geq5.
\]
由\cref{thm:canonical-reduction}，第二层构成满足条件的 $(2,3)$ blocker pair；
第三层交换两侧后也满足同一条件。\cref{thm:m23}因此给出
$|\mathcal A_2|,|\mathcal A_3|\geq9$。所以
\[
  |\family|\geq1+5+9+9+5+1=30.
\]
上下界相等。
这里的 $3\cdot2^{30}$ 只是该论证使用的一个方便阈值见证，并不声称它是最优
稳定化阈值。
\end{proof}

\section{七层情形的全局下界}
\label{sec:sat7-lower}

七层情形所需的技术 blocker 估计见\cref{app:m24}。

\begin{theorem}
\label{thm:m24}
精确局部 blocker 参数为 $m(2,4)=12$。
\end{theorem}

\begin{proof}
由\cref{prop:m24-product}，每个满足条件的 pair 都有 $ab\geq32$，其中
$a,b$ 为两侧大小。因此
\[
  (a+b)^2=(a-b)^2+4ab\geq128>11^2,
\]
故 $a+b\geq12$。

取两个顶点不交四环的八条边组成 $\mathcal S$。每个四环的 blocker 恰由两个
二分部组成，所以 $\blocker(\mathcal S)$ 是从每个环各取一个二分部后得到的
四个并集。反过来，这四个并集的任意横截集都包含某个环的一条边。因此两族互为
blocker，成员大小下界分别为二和四，总基数为 $8+4=12$。
\end{proof}

\begin{lemma}
\label{lem:middle-nine}
极小饱和 7-Sperner 族的 canonical 分解中，
$|\mathcal A_3|\geq9$。
\end{lemma}

\begin{proof}
使用\cref{thm:canonical-reduction}的记号。$\mathcal S_3$ 与
$\mathcal C_3$ 的每个成员大小至少为三。从 $U$ 的所有子集中等概率选取 $R$。
Canonical 覆盖说明，至少一个事件 $S\subseteq R$ 或
$R\cap C=\varnothing$ 发生。每个事件的概率至多为 $2^{-3}$。并集界给出
$|\mathcal A_3|\geq8$。

若等号成立，则每个事件的概率都为 $1/8$，且这些事件两两不交。$R=U$ 属于
每个小侧事件，$R=\varnothing$ 属于每个大迹补事件。两侧都非空，所以等号会
迫使每侧各仅有一个事件，这与总数为八矛盾。因此
$|\mathcal A_3|\geq9$。
\end{proof}

\begin{theorem}
\label{thm:sat7-lower}
最终稳定的饱和 7-Sperner 数满足
\[
  \sat(7)\geq47.
\]
\end{theorem}

\begin{proof}
选取满足 $n\geq3\cdot2^{55}$ 的稳定底集。\cref{thm:size55}的构造说明极小
族至多有 55 个成员，所以\cref{thm:canonical-reduction}的两个数值条件均成立。
该定理与私有见证给出
\[
  (|\mathcal A_0|,|\mathcal A_1|,|\mathcal A_5|,
  |\mathcal A_6|)\geq(1,6,6,1).
\]
典范迹 blocker pair 与\cref{thm:m24}给出
$|\mathcal A_2|,|\mathcal A_4|\geq12$；
\cref{lem:middle-nine}给出 $|\mathcal A_3|\geq9$。七层求和可得
\[
  \sat(7)\geq1+6+12+9+12+6+1=47.
\]
\end{proof}

\section{含 55 个集合的构造}
\label{sec:construction}

令 $C=[8]=\{1,\ldots,8\}$，令 $H$ 与 $C$ 不交，并假设 $|H|>2$。
我们把核心子集的元素连接起来作为缩写；例如，$145$ 表示
$\{1,4,5\}$，而 $1346H$ 表示 $\{1,3,4,6\}\cup H$。符号 $-$ 表示
空列表。

对\cref{tab:layers}的每一行，令 $\smallcores_i$ 为列出的小核心，令
$\largecores_i$ 为列出的大核心。定义
\[
  \mathcal A_i
   =\smallcores_i\cup\{L\cup H:L\in\largecores_i\},
  \qquad 0\leq i\leq6.
\]

\begin{table}[htbp]
\centering
\small
\setlength{\tabcolsep}{4pt}
\renewcommand{\arraystretch}{1.18}
\begin{tabular}{c p{0.35\textwidth} p{0.48\textwidth}}
\toprule
$i$ & $\smallcores_i$ & $\largecores_i$ \\
\midrule
0 & $\varnothing$ & -- \\
1 & $1$, $2$, $3$, $6$, $7$ & $458$ \\
2 & $24$, $15$, $35$, $26$, $47$, $67$, $18$, $38$
  & $1346$, $1237$, $4568$, $2578$ \\
3 & $124$, $234$, $145$, $345$, $156$, $356$, $267$, $178$, $378$
  & $12357$, $13467$, $12358$, $12368$, $13468$, $24568$, $24578$, $45678$ \\
4 & $2345$, $1247$, $3568$, $1678$
  & $123567$, $134567$, $123468$, $124568$, $123578$, $134578$, $234678$, $245678$ \\
5 & $23457$
  & $1234568$, $1234678$, $1235678$, $1245678$, $1345678$ \\
6 & -- & $12345678$ \\
\bottomrule
\end{tabular}
\caption{七层的核心描述。右栏中的每个大核心都代表它与 $H$ 的并。}
\label{tab:layers}
\end{table}

\begin{proposition}
\label{prop:layer-certificate}
对每个 $i\in\{0,\ldots,6\}$，族 $\mathcal A_i$ 都是饱和反链。
\end{proposition}

\begin{proof}
在\cref{tab:layers}每个非空的小核心列中，所有集合都有相同的基数，所以
$\smallcores_i$ 是反链。对每一行，大核心列恰好由那些避开小核心列全部成员的
包含极大核心子集组成。这一陈述是一个有限证书：向列出的大核心中加入任意一个
缺失的核心元素，都会产生一个列出的小核心；而每个不含小核心的核心集合，都
包含于某个列出的大核心中。后一覆盖蕴含在每一行中只有 $2^8$ 个核心掩码；
补充验证器检查了全部七行中的这两个蕴含。两个端点行分别表示：没有集合能避开
$\varnothing$；当小核心族为空时，$C$ 是唯一的极大避集。结论由
\cref{lem:maximal-free}立即得到。
\end{proof}

为完整起见，附录~\ref{app:layering}为第零层以上的每个成员记录了一个前驱。
所有列出的包含关系都是严格的，因此这七层是分层的。\cref{tab:layers}中的
55 个模板在不同层之间两两不同。

\begin{theorem}
\label{thm:size55}
对每个与 $[8]$ 不交且满足 $|H|>2$ 的有限集 $H$，族
\[
  \family=\bigcup_{i=0}^{6}\mathcal A_i
\]
是一个基数为 55 的饱和 7-Sperner 族。
\end{theorem}

\begin{proof}
由\cref{prop:layer-certificate}，全部七层都是饱和反链。
附录~\ref{app:layering}中的前驱证书证明了分层性，而直接检查
\cref{tab:layers}可知各层两两不交。因此，\cref{thm:layered}说明它们的并
是饱和 7-Sperner 族。各层大小为
\[
  1, 6, 12, 17, 12, 6, 1,
\]
总和为 55。
\end{proof}

\begin{proposition}
\label{prop:endpoints-atom}
族 $\family$ 同时含有 $\varnothing$ 和 $C\cup H$，并且 $H$ 是
$\family$ 的一个齐次原子。相对于这一原子，$\family$ 有 28 个小成员和
27 个大成员。
\end{proposition}

\begin{proof}
两个端点分别是 $\mathcal A_0$ 和 $\mathcal A_6$ 的唯一成员。每个列出的模板
要么包含整个 $H$，要么与 $H$ 不交。为了验证极大性，注意每个核心点都出现在
某个小成员中：点 $1,2,3,6,7$ 出现在 $\mathcal A_1$ 的单点成员中，而
$4,5,8$ 分别出现在 $24$、$15$ 和 $18$ 中。因此，对每个 $c\in C$，
$\family$ 中某个成员与 $H\cup\{c\}$ 的交是一个非空真子集。于是，$H$ 的
任何严格超集都不可能具有原子性质。由于 $|H|>2$，该原子是齐次的。最后，
对\cref{tab:layers}两列计数可得，小模板数为
$1+5+8+9+4+1=28$，大模板数为
$1+4+8+8+5+1=27$。
\end{proof}

\par\medskip

\begin{corollary}
\label{cor:main-bound}
对每个 $n\geq11$，都有 $\sat(n,7)\leq55$。因此 $\sat(7)\leq55$。
\end{corollary}

\begin{proof}
取 $|H|=n-8$。此时 $|H|>2$，故\cref{thm:size55}适用于一个 $n$ 元底集。
关于稳定值的不等式由 $\sat(7)$ 的定义得到。
\end{proof}

\section{排除 54 元族}
\label{sec:exact-seven}

下面改进下界。设 $\mathcal A_0,\ldots,\mathcal A_6$ 是稳定底集上某个极小族
的 canonical 层，并取一个齐次原子。对每个内部层，记 $\mathcal S_i$ 为其小迹，
记 $\mathcal C_i$ 为大迹的补族，定义同\cref{thm:canonical-reduction}。
《补充证明与验证》包含较长的任意有限归约，以及一条从 blocker pair 到有限分类
的完整示例。

\begin{theorem}[中间 blocker 下界]
\label{thm:m33}
精确局部参数为 $m(3,3)=14$。
\end{theorem}

\begin{proof}
补充材料 S2 节给出任意有限下界证明。私有见证把每行大小限制在对侧行数以内；
最初的度数与残余归约先给出下界 13。若取等，不计 blocker 对偶时只剩
$(4,9)$、$(5,8)$、$(6,7)$。残余 blocker 排除高点度；当每点度至多为二时，
实际点形成行标签上的自环多重图，并保留所有平行实际点。包含极小支撑核具有
有界顶点度，因此支撑与重数域有限。精确分类排除三个拆分，得到下界 14。

七条 Fano 直线构成自 blocker clutter，达到总数 14，所以下界可取到。
\end{proof}

\begin{theorem}[中间等号与单松弛分类]
\label{thm:m33-classification-main}
总基数为 14 的任意合格 pair，在删除非活跃点并重标活跃底集后，都是 Fano
直线 clutter 与自身组成的 pair。不存在总基数为 15 的合格 pair。
\end{theorem}

\begin{proof}
补充材料 S3 节给出等号与单松弛归约。残余与支撑论证排除非对称等号拆分；中央
$(7,7)$ 拆分具有 Fano 关联模式，blocker 对合固定另一侧。加入一个松弛行后重复
归约，可排除全部总数 15 拆分。S10 节的商与反向展开结果覆盖非活跃点、平行点、
重复关联模式与非一致行。因此分类适用于原始的任意有限底集。
\end{proof}

\begin{proposition}[Canonical 轮廓归约]
\label{prop:canonical-profiles}
不存在总大小为 52 或 53 的 canonical 族。如果 canonical 族的总大小为 54，
其层轮廓只能是
\[
  P_{\mathrm F}=(1,6,13,14,13,6,1),\qquad
  P_{16}=(1,6,12,16,12,6,1).
\]
第一个轮廓的中间 blocker 对是 Fano 直线 clutter。第二个轮廓的中间 blocker
对总基数为 16。
\end{proposition}

\begin{proof}
有限典范归约、私有见证、\cref{thm:m24}与\cref{thm:m33}给出
\[
  (|\mathcal A_0|,\ldots,|\mathcal A_6|)
  \geq(1,6,12,14,12,6,1).
  \tag{2}
\]
由\cref{thm:m33-classification-main}，式 (2) 取等时，中间 blocker 对只能是
Fano 平面。它无法与下方或上方的 12 元等号相邻层兼容，所以总数 52 不可能。
若仅增加一个成员，则仍落入 Fano 排除，或使中间 blocker 对的总数为 15；后者
由\cref{thm:m33-classification-main}排除。所以总数 53 也不可能。

对总数 54，在式 (2) 的五个内部层上引入非负松弛量，其和为二。若中间层有
14 个成员，则它是 Fano 对，并迫使两个相邻层都严格大于等号。15 元中间对不
存在。若中间层超过 15，则它恰有 16 个成员，其他各层全部取等。由此仅剩所列
两个轮廓。
\end{proof}

为证明精确排除，令
\[
\begin{aligned}
  \mathcal F&=\mathcal S_2,&\mathcal G&=\mathcal C_2,\\
  \mathcal H&=\mathcal S_3,&\mathcal K&=\mathcal C_3,\\
  \mathcal P&=\mathcal S_4,&\mathcal Q&=\mathcal C_4.
\end{aligned}
\]
相邻接口说明：$\mathcal G$ 的每个成员都严格包含 $\mathcal K$ 的某个成员，且
$\mathcal F\cap\mathcal H=\varnothing$。对偶地，$\mathcal P$ 的每个成员都
严格包含 $\mathcal H$ 的某个成员，且
$\mathcal K\cap\mathcal Q=\varnothing$。

\begin{proposition}[54 元排除]
\label{prop:exclude54}
\cref{prop:canonical-profiles}中的两个轮廓在任意有限迹底集上都无法实现。
\end{proposition}

\begin{proof}
先考虑 $P_{\mathrm F}$。每条 Fano 直线都严格包含 $\mathcal F$ 中一个不同的
二元组。不计 blocker 对偶时，相邻拆分为 $(7,6)$、$(8,5)$ 或 $(9,4)$。
在第一个拆分中，$\mathcal F$ 从每条 Fano 直线的三个二元组中选择一个。穷尽
检查全部 $3^7=2187$ 个选择后，没有任何选择恰有六个大小至少为四的极小点覆盖。
在第二个拆分中，写成 $\mathcal F=E(Q)\cup\{R\}$。迹 $R\cap X$ 是 $Q$ 中
的独立集，而 $R$ 的外部点构成一个关联类。私有见证界说明其精确重数至多为五。
这样，每个有限实例都归结为 249576 个三元组 $(Q,R\cap X,t)$ 之一。没有合格
三元组恰有五个大小至少为四的 blocker。最后一个拆分中，$\mathcal G$ 的四行
大小至少为四，且两两相交。残余迹论证给出至少十个 blocker，与
$|\mathcal F|=9$ 矛盾。因此 Fano 轮廓及其上方对偶均不可能。

再考虑 $P_{16}$。对有限底集均匀随机取子集。两个中间 blocker 侧对应事件
$D\subseteq R$ 与 $E\subseteq U\setminus R$ 构成穷尽覆盖。相反方向的两个
事件不交，因为每个 $D$ 都与每个 $E$ 相交；同一侧事件未必不交。若 16 行大小
都至少为四，则并合界必须取等，并迫使全部事件两两不交。私有见证说明每侧至少
有三行，而同一侧任意两个事件以正概率相交，矛盾。因此至少有一条中间行大小为三。

不计交换 blocker 两侧时，中间拆分为 $(5,11)$、$(6,10)$、$(7,9)$ 或
$(8,8)$。下表记录全部归约。所有重数都按实际点计数，所以平行支撑和重复关联
模式仍保持为不同的点。
\begin{table}[htbp]
\centering
\small
\setlength{\tabcolsep}{5pt}
\renewcommand{\arraystretch}{1.16}
\begin{tabular}{c p{0.41\textwidth} p{0.38\textwidth}}
\toprule
拆分 & 任意有限归约 & 穷尽余项 \\
\midrule
$(5,11)$ & 一个三度点给出九个含该点的 blocker，以及两个避开该点的残余
blocker。 & 若处处取等，残余将成为两个大小至少为三的不交行的 blocker；这会
从仅有五行的残余中产生至少九行。 \\
$(6,10)$ & 度至少为四时可得十一个 blocker。度为三时留下三个精确三行核；
度至多为二时得到一个实际自环多重图。 & 非解析核有 95 个重数向量，没有兼容
扩展。每个低度核至少有 15 个 blocker。 \\
$(8,8)$ & 度数归约后最大度为三。一条三元行留下一个支撑秩至多为三的五行核。
& 共有 1890 个重数向量，分成 24 类；2514 个可能扩展中没有合格实例。 \\
$(7,9)$ & 两侧最大度分别为三和四。若九行侧有三元行，则留下四行核；若三元行
仅在七行侧，则必须耦合两个三行残余。 & 第一支有 136 个向量和 138161 个扩展，
没有合格实例。第二支的 90898 个耦合全部至少有四个二元 blocker，与上界三
矛盾。 \\
\bottomrule
\end{tabular}
\caption{16 元轮廓中全部中间拆分的穷尽归约。}
\label{tab:exact-splits}
\end{table}


行数、私有见证、点度界和支撑极小化给出\cref{tab:exact-splits}中的有限域。
补充材料 S5--S10 节证明支撑商可以反向展开为任意有限 blocker pair。独立的
C++ 与 Python 程序检查所得有限核，Lean 的可靠性与实现引理再把结果连接到原始
clutter。因此四个拆分及其 blocker 对偶均不可能。被强制存在的三元行无处可放，
从而排除 $P_{16}$。
\end{proof}

\begin{theorem}
\label{thm:sat7}
最终稳定的饱和 7-Sperner 数为
\[
  \sat(7)=55.
\]
\end{theorem}

\begin{proof}
在满足 $n\geq3\cdot2^{55}$ 的稳定底集上选择一个极小族。55 元构造把其大小
限制为 55，故\cref{thm:canonical-reduction}适用。52 元与 53 元排除给出
$\sat(7)\geq54$，而\cref{prop:exclude54}排除等号。因此
$\sat(7)\geq55$。反向不等式由\cref{cor:main-bound}给出。
\end{proof}

\section{复合推论}
\label{sec:composition}

由\cref{prop:endpoints-atom}，\cref{thm:size55}中的族满足
Morrison--Noel--Scott 复合引理~\cite[Lemma~18]{mns2014}所要求的端点条件
和齐次原子条件。该引理在不交底集上把两个饱和系统的“小--小”部分和
“大--大”部分分别配对。若两个因子的小成员数与大成员数分别为
$(a_1,b_1)$ 和 $(a_2,b_2)$，则复合后的数量为 $(a_1a_2,b_1b_2)$，
饱和参数为 $k_1+k_2-2$；同时，端点条件和齐次原子条件仍被保留。

因此，将这个 55 元因子迭代 $j\geq1$ 次，所得参数为 $5j+2$，小成员数与
大成员数分别为 $(28^j,27^j)$。对余项因子，取一个 $s$ 元集合 $Y$ 和另一个
齐次原子 $K$，并使用
\[
  \powerset(Y)\ \cup\ \{S\cup K:S\subseteq Y\}.
\]
这是标准的饱和 $(s+2)$-Sperner 系统：其小部分和大部分的大小均为 $2^s$，
并且它含有两个端点。将它与上述迭代因子复合，会使参数增加
$(s+2)-2=s$，并分别得到大小为 $2^s28^j$ 和 $2^s27^j$ 的小部分与
大部分。我们取 $0\leq s\leq4$，以覆盖通常的五个剩余类。

\begin{corollary}
\label{cor:discrete-bound}
对 $j\geq1$ 和 $0\leq s\leq4$，有
\[
  \sat(5j+2+s)\leq 2^s\bigl(28^j+27^j\bigr).
\]
\end{corollary}

作为比较，对称的 56 元因子给出 $2^{s+1}28^j$。对每个 $j\geq1$，
不等式
\[
  2^s(28^j+27^j)<2^{s+1}28^j
\]
都是严格的。主导指数增长率没有改变，因为较大部分的大小仍为 28。

\section{七点核心模板类中的最优性}
\label{sec:class-optimality}

55 元构造使用了八元素核心。我们还检验了相应的七点核心模板类能否包含一个
55 元族。下面的结果是计算机辅助的，其适用范围有意取得比
\cref{thm:size55}更窄。

\begin{proposition}[证书支持的类内最优性]
\label{prop:class-optimum}
设 $|C|=7$ 且 $|H|>2$。在所有满足下述条件的有序七元组中，总基数的最小值
为 56：七个族两两不交、构成分层的饱和反链，并且每个成员都形如 $A$ 或
$A\cup H$，其中 $A\subseteq C$。
\end{proposition}

\begin{proof}[计算机辅助证明]
对层 $i$、原子使用位 $b$ 和核心掩码 $m$，引入布尔变量 $x_{i,b,m}$，表示
是否选择相应模板。CNF 包含：同一层内可比模板对的排除子句；对三种齐次块
特征和每个核心掩码各设置一个覆盖子句；不同层之间的两两排除子句；用于分层的
前驱蕴含；以及一个顺序计数器~\cite{sinz2005}，把所选模板总数限制为至多 55。

三种特征分别记录任意测试集合与 $H$ 的交是 $\varnothing$、非空真子集，还是
整个 $H$。一个集合与模板之间的包含关系只取决于这一特征和核心掩码。因此，
每个大小至多为 55 的数学候选都能扩展为 CNF 的一个满足赋值。反过来，一个
满足赋值中被选中的变量可以解码为命题所述的模板条件；另外还包含独立语义检查
作为对照。

这个确定性实例包含 100,297 个变量和 252,427 个子句。Kissat 4.0.4 返回
UNSAT，并生成一个大小为 253,731,315 字节的 DRAT 证明。
DRAT-trim~\cite{wetzler2014} 独立检查该证明，并返回
\texttt{s VERIFIED}。所以，不存在大小至多为 55 的七点核心候选。
Martin--Veldt 的 56 元构造给出了与之匹配的上界。
\end{proof}

该编码既没有施加补集对称性，也没有预先指定各层大小，并且不要求 Fano 平面
模式。尽管如此，\cref{prop:class-optimum}并不覆盖更大的核心、非模板族或不含
公共块的系统，因此它没有给出 $\sat(7)$ 的全局下界。

\section{形式化与计算验证}
\label{sec:verification}

构造与下界以不同方式使用计算。一个符号程序对每一层检查
$3\cdot2^8=768$ 种核心/块特征。另有三项检查使用十一点实现的完整幂集，其中
两项以独立算法检验 55 元并集。四项检查均接受\cref{tab:layers}中的族，并能
定位早期 54 元构造失效的层。

Lean 4~\cite{lean4} 与 Mathlib~\cite{mathlib2020} 验证了十一点实现。形式化
定理证明该族包含 55 个成员，是 7-Sperner 族且饱和。饱和性表覆盖全部
$2^{11}=2048$ 个掩码。这条有限定理复核具体实现；\cref{thm:size55}则直接证明
任意 $|H|>2$ 的情形。

主要 Lean 开发直接形式化最终稳定值。相应谓词要求存在阈值 $N$。对每个
$n\geq N$，既存在大小为 $s$ 的饱和 $k$-Sperner 族，任意这样的族也至少有
$s$ 个成员。两个最终声明为
\begin{quote}\small\ttfamily
sat\_six\_eq\_thirty : IsStableSaturationNumber 6 30\\
sat\_seven\_eq\_fifty\_five : IsStableSaturationNumber 7 55.
\end{quote}

六层证明把 30 元构造、任意有限定理 $m(2,3)=9$、齐次原子、典范分解及分层
求和连接起来。七层证明结合 55 元构造、典范轮廓归约、Fano 相邻层排除，以及
总数 16 的四种中间拆分。形式化还包含从任意有限 clutter 到支撑商的归约和
反向实现映射。

最长的有限分类没有在正文逐项展开。对应的 Lean 声明证明支撑归约、实现映射和
各分支排除；独立的 C++ 与 Python 程序检查同一批有限核。Supplement 给出数学
引理、程序与 Lean 声明之间的精确对应。归档工件固定 Lean 与 Mathlib 版本，
并提供完整构建方法。

\section{讨论与开放问题}

构造与下界回答了 Martin--Veldt 的问题：七层最终稳定值为 55。把对称划分
$28/28$ 替换为 $28/27$ 不改变主导渐近指数，却改进了
\cref{cor:discrete-bound}中的每个有限上界。

七点核心类的结果解释了原构造为何无法在同一模板内继续压缩。其公共块模板类的
最小值为 56，新构造则使用第八个核心点。其他 55 元构造未必属于这一模板，
也未必能用齐次块描述。

同一 blocker 观点补上了已知 30 元六层构造的下界。它把稳定下界问题拆成相邻
典范层之间的局部问题。七层时，这些局部问题仍可精确分类。下一个未知值是
$\sat(8)$，需要更强的 blocker 估计或新的结构。

\par\medskip\noindent\textbf{补充材料。}\quad
随稿单独提交的《补充证明与验证》包含较长的 blocker 归约、有限域证明、计算
检查，以及它们与 Lean 形式化的对应关系。

\section*{可复现性与声明}
\addcontentsline{toc}{section}{可复现性与声明}

\par\medskip\noindent\textbf{数据与代码可得性。}\quad
本研究不使用经验数据。配套工件包含证明、搜索程序、构造验证器、证书、Lean
源代码及构建方法。该工件可从 GitHub 发布
\href{https://github.com/lailai0916/saturated-sperner-6-7/releases/tag/v1.1.2}
{v1.1.2} 获取；Zenodo 归档 DOI 为
\href{https://doi.org/10.5281/zenodo.21770438}{10.5281/zenodo.21770438}。

\par\medskip\noindent\textbf{生成式 AI 与 AI 辅助技术声明。}\quad
OpenAI Codex 辅助了文献检索规划、证明推导、代码起草、验证和语言编辑。作者
核对了数学结论、代码与参考文献，并对最终论文负责。

\par\medskip\noindent\textbf{资助。}\quad
本研究未获得公共、商业或非营利部门资助机构的任何专项资助。

\par\medskip\noindent\textbf{利益冲突声明。}\quad
作者声明，不存在可能影响本文工作的已知竞争性经济利益或个人关系。

\appendix

\section{局部 blocker 参数的乘积不等式}
\label{app:m24}

\begin{proposition}
\label{prop:m24-product}
设 $\mathcal S,\mathcal C$ 是有限底集上的 mutual blockers，
$\mathcal S$ 的每个成员大小至少为二，$\mathcal C$ 的每个成员大小至少为四。
则
\[
  |\mathcal S||\mathcal C|\geq32.
\]
\end{proposition}

\begin{proof}
记 $a=|\mathcal S|$、$b=|\mathcal C|$。私有见证引理给出
$a\geq4$ 与 $b\geq2$。下面证明六个下界：
\[
\begin{array}{cc@{\qquad\qquad}cc}
\toprule
b&a\text{ 的下界}&a&b\text{ 的下界}\\
\midrule
2&16&4&16\\
3&12&5&8\\
4&8&6&6\\
\bottomrule
\end{array}
\tag{2}
\]

若 $b=2$，$\mathcal C$ 的两个成员不交，逐边取点给出至少
$4^2=16$ 个 blocker。若 $b=3$，使用\cref{thm:m23}证明中的六个 incidence
类。式 (1) 仍成立，但行不等式变为
$\alpha_i+\pi_j+\pi_k\geq4$。若某个 $\pi_i\geq3$，相关项至少贡献
$3\cdot4=12$。否则令
$0\leq\pi_1\leq\pi_2\leq\pi_3\leq2$，并把各 $\alpha_i$ 取到最小。
十种情形依次给出
\[
  64,40,24,25,16,12,17,13,12,12,
\]
所以 $a\geq12$。若 $a=4$，$\mathcal S$ 的每个横截集大小至少为四。任意
两行若相交，则该交点与其余两行各一点会构成至多三点的横截集。因此四行两两
不交，逐行取点给出 $b\geq2^4=16$。

下面证明 $b=4$ 的界。写
$\mathcal C=\{C_1,C_2,C_3,C_4\}$，并取大小最小的成员 $E=C_1$，记
$n=|E|\geq4$。对 $j=2,3,4$，令
\[
  O_j=E\setminus C_j,\qquad R_j=C_j\setminus E,\qquad
  o_j=|O_j|,\quad q_j=|R_j|.
\]
Clutter 性说明所有 $O_j,R_j$ 非空；$E$ 的最小性给出 $q_j\geq o_j$；没有
单点 blocker 给出 $O_2\cup O_3\cup O_4=E$。对 $x\in E$，定义
\[
  \mathcal J_x=\{R_j:x\in O_j\}.
\]
由私有见证，映射 $R\mapsto\{x\}\cup R$ 把
$\blocker(\mathcal J_x)$ 双射到那些与 $E$ 的交为 $\{x\}$ 的
$\blocker(\mathcal C)$ 成员。因此，若
$w_x=\min_{j:x\in O_j}o_j$，则
\[
  |\blocker(\mathcal C)|\geq\sum_{x\in E}w_x.
\tag{3}
\]
把每个 $x$ 分配给一个取到 $w_x$ 的下标，并记三个类大小为 $t_j$。由于
$t_j\leq o_j$ 且 $\sum t_j=n$，有
\[
  \sum_{x\in E}w_x=\sum_{j=2}^4t_j o_j
  \geq\sum_{j=2}^4t_j^2
  \geq\left\lceil\frac{n^2}{3}\right\rceil.
\tag{4}
\]
若 $n\geq5$，右端至少为九。以下设 $n=4$。式 (3) 的右端至少为六。若它
等于七，四个非空迹
$E,E\setminus O_2,E\setminus O_3,E\setminus O_4$ 有一个大小至少为二的
极小横截集；它给出 singleton-intersection 类之外的另一个 blocker。若右端
等于六，式 (4) 取等迫使三个 $O_j$ 按 $2,1,1$ 划分 $E$。两点块与两个单点
块之间的四个 cross-pair，再加上两个单点组成的 pair，给出五个额外 blocker。
所以所有情形都有 $|\blocker(\mathcal C)|\geq8$。

若 $a=5$，$\mathcal S$ 的每个横截集大小至少为四。每个点至多属于两行。
行交图的匹配数至多为一，否则两条不交边和第五行的一点会给出三点横截集。
这种图只可能是空图、星或三角形，并允许添加孤立点。空图给出 $2^5$ 个 blocker。
星中固定一个中心--叶交点，再从其他三条两两不交的行各取一点。三角形中固定
一条边上的交点，再从第三条三角形行和各孤立行取点；incidence degree 至多为二，
所以固定点仍有私有见证。两个非空情形都至少给出 $2^3=8$ 个 blocker。

最后设 $a=6$。若一点属于三行，则其余三行两两不交；固定该点并从其余三行
各取一点，得到至少八个 blocker。否则，按每个实际点属于一行或两行，把它编码
为六个行标签上的 loop 或 edge；平行实际点保持区分。每个顶点的 incidence
degree 至少为二，非 loop support 图的匹配数至多为二。

若匹配数为零，loop choices 给出至少 $2^6$ 个极小 cover。若匹配数为一，
nonloop support 由一个星形图或一个三角形及若干孤立顶点组成。星形情形中，固定一个
center--leaf 实际点，再从其余四个顶点各取一个 incident point。这四个顶点间
没有 nonloop support，所以所选点两两不同且各有 private vertex，至少得到
$2^4$ 个 cover。三角形情形中，固定两条相邻边上的实际点；其余三个顶点在
nonloop support 中孤立，且各有至少两个 loop choices，故至少得到 $2^3$ 个
cover。

以下设匹配数为二。固定两个不交 support $E_1,E_2$，其 multiplicity 为
$m_1,m_2$；记未匹配顶点为 $u,v$。从每个 $E_i$ 取一个实际点，再分别取一个
与 $u,v$ incident 的点，可得到
\[
  m_1m_2d(u)d(v)
\]
个不同的四点 cover。由于不存在三点横截集，它们都是极小的。除非
$m_1=m_2=1$ 且 $d(u)=d(v)=2$，该数至少为六。

在这个例外中，第二个实际二边匹配会产生至少四个同类 cover。两个生成族的交
至多为二。若两个匹配不共享实际点，则共同 cover 只能是二者的四点并。若二者
共享一个实际点，则余下两个 support 不可能不交；若它们交于一个顶点，第四点
必须覆盖一个 degree 为二的固定未匹配顶点；若它们平行，则不存在共同的极小
cover。因此两族的并至少含六个 cover。

若固定匹配唯一，则 $u,v$ 的 incidence 都是 loop。令 $e_i$ 为 $E_i$ 上唯一的
实际点。再写 $E_1=\{a,b\}$、$E_2=\{c,d\}$，并在 $a,b$ 处分别选取
$r_a\neq e_1$、$r_b\neq e_1$。唯一性迫使 $r_a\neq r_b$，且二者的
support 都位于四个已匹配顶点中，所以 $\{r_a,r_b,e_2\}$ 覆盖这四个顶点。
从中取一个避开 $e_1$ 的包含极小子 cover $K$。把 $K$ 分别与 $u,v$ 处的两个
loop choice 组合，得到四个避开 $e_1$ 的极小 cover；它们与原来包含 $e_1$
的四个 cover 不交。因此 $b\geq6$。

式 (2) 穷尽全部可能：
\[
\begin{array}{cc}
\toprule
\text{情形}&ab\\
\midrule
b=2&\geq32\\
b=3&\geq36\\
b=4&\geq32\\
b\geq5,\ a=4&\geq64\\
b\geq5,\ a=5&\geq40\\
b\geq5,\ a=6&\geq36\\
b\geq5,\ a\geq7&\geq35\\
\bottomrule
\end{array}
\]
命题得证。
\end{proof}

\section{分层前驱证书}
\label{app:layering}

在下表中，$B\leftarrow A$ 表示 $A$ 属于前一层且 $A\subsetneq B$。
后缀 $H$ 表示与整个齐次块取并。

\begin{description}[leftmargin=1.8em,style=nextline]
\item[$\mathcal A_1\leftarrow\mathcal A_0$]
$1\leftarrow\varnothing$;
$2\leftarrow\varnothing$;
$3\leftarrow\varnothing$;
$6\leftarrow\varnothing$;
$7\leftarrow\varnothing$;
$458H\leftarrow\varnothing$.

\item[$\mathcal A_2\leftarrow\mathcal A_1$]
$24\leftarrow2$; $15\leftarrow1$; $35\leftarrow3$;
$26\leftarrow2$; $47\leftarrow7$; $67\leftarrow6$;
$18\leftarrow1$; $38\leftarrow3$; $1346H\leftarrow1$;
$1237H\leftarrow1$; $4568H\leftarrow6$; $2578H\leftarrow2$.

\item[$\mathcal A_3\leftarrow\mathcal A_2$]
$124\leftarrow24$; $234\leftarrow24$; $145\leftarrow15$;
$345\leftarrow35$; $156\leftarrow15$; $356\leftarrow35$;
$267\leftarrow26$; $178\leftarrow18$; $378\leftarrow38$;
$12357H\leftarrow15$; $13467H\leftarrow47$;
$12358H\leftarrow15$; $12368H\leftarrow26$;
$13468H\leftarrow18$; $24568H\leftarrow24$;
$24578H\leftarrow24$; $45678H\leftarrow47$.

\item[$\mathcal A_4\leftarrow\mathcal A_3$]
$2345\leftarrow234$; $1247\leftarrow124$;
$3568\leftarrow356$; $1678\leftarrow178$;
$123567H\leftarrow156$; $134567H\leftarrow145$;
$123468H\leftarrow124$; $124568H\leftarrow124$;
$123578H\leftarrow178$; $134578H\leftarrow145$;
$234678H\leftarrow234$; $245678H\leftarrow267$.

\item[$\mathcal A_5\leftarrow\mathcal A_4$]
$23457\leftarrow2345$; $1234568H\leftarrow2345$;
$1234678H\leftarrow1247$; $1235678H\leftarrow3568$;
$1245678H\leftarrow1247$; $1345678H\leftarrow3568$.

\item[$\mathcal A_6\leftarrow\mathcal A_5$]
$12345678H\leftarrow23457$.
\end{description}

\phantomsection
\addcontentsline{toc}{section}{参考文献}
{\small
\bibliographystyle{plain}
\bibliography{references}
}

\end{document}
